\section{Related Work}
\label{sec:related-work}

\paragraph{External agent memory and long-term memory evaluation.}
Language agents externalize experience in several forms: Generative Agents stores
observations and higher-level reflections \citep{park2023generative}, Reflexion
reuses verbal feedback from prior trials \citep{shinn2023reflexion}, and ExpeL
extracts transferable insights from successful and failed trajectories
\citep{zhao2023expel}. Long-term memory benchmarks complement these systems by
testing multi-session recall, event summarization, temporal reasoning, knowledge
updates, and abstention \citep{maharana2024locomo,wu2025longmemeval}. Together,
this literature establishes that persistent experience can guide later behavior
and that retrieval and updating require explicit evaluation. Our setting examines
a complementary operational question: whether a previously validated strategy
revision should still be reused after public demand or supply conditions have
changed and outcome feedback arrives only after a delayed window. ShiftMem
therefore evaluates memory through downstream inventory cost and recovery while
making suspension, revalidation, and reactivation auditable, rather than treating
recall accuracy alone as evidence of operational applicability.

\paragraph{Concept drift and nonstationary inventory control.}
Concept-drift detection and adaptive inventory control address different parts of
decision making under change. ADWIN statistically compares subwindows
\citep{bifet2007adwin}, cumulative-sum procedures detect sustained process-level
departures \citep{page1954continuous}, and broader drift research separates
detection from the adaptation strategy that follows \citep{gama2014survey}.
Inventory research instead develops domain-specific control under changing or
partially observed demand \citep{treharne2002adaptive,bayraktar2010inventory},
including lost-sales settings with censored demand and replenishment delay
\citep{zipkin2008old} and online control under unknown changing demand and lead
times \citep{amiri2026nonstationary}. ShiftMem neither treats a detector as a
validity oracle nor claims a new optimal inventory policy. A Page--Hinkley signal
computed from public observations triggers an additional review and, when
detector and experience variable names match, can return an active memory to
probation. Stored lead-time conditions gate retrieval, while delayed outcomes
update lifecycle confidence and state. Holding the bounded deterministic
controller fixed allows an end-to-end comparison of the two implemented memory
systems; it does not identify the effect of an individual lifecycle component.

\paragraph{Operational LLM reliability and reproducible evaluation.}
LLM inventory agents can coordinate or propose supply-chain decisions
\citep{quan2024invagent}, but their behavior varies across models and uncertain
inventory settings \citep{zhao2025aimbench}. ShiftMem consequently separates
language-model reasoning from routine numerical execution: the LLM reviews only
three bounded strategy parameters, cannot issue daily orders or observe hidden
regime and Oracle state, and falls back to the previous valid strategy after one
failed schema-repair attempt. This reliability path remains part of the measured
system because format restrictions can affect model performance
\citep{tam2024letmespeak}, while grammar-constrained decoding treats output
validity as a distinct systems problem \citep{geng2023grammarconstrained}. The
evaluation also follows work on pre-specification, reproducibility, and transparent
reporting \citep{nosek2018preregistration,pineau2021reproducibility,
dodge2019showyourwork}. We use the narrower term \emph{pre-Test-frozen}, not
formal preregistration, and report the unsupported primary hypothesis with its
effect estimate, uncertainty, heterogeneity, and post-hoc sensitivity analyses
rather than reducing the result to a binary $p$-value conclusion
\citep{wasserstein2016asa}. The resulting contribution is an empirical test of a
conditional memory lifecycle within a bounded operational review system, not a
claim of universal memory superiority.
